\documentclass[10pt,a4paper]{amsart}
\usepackage[margin=19mm]{geometry}
\usepackage{amsmath,amssymb,mathtools}
\usepackage[scr=rsfs,cal=euler]{mathalfa}
\usepackage{xfrac}
\usepackage[unicode,hidelinks,bookmarks=false]{hyperref}
\newcommand{\ie}{i.e.\ }
\newcommand{\cf}{cf.\ }
\newcommand{\resp}{respectively\ }
\newcommand{\Subsection}{\subsection}
\providecommand{\nobreakdash}{\nobreak\hbox{-}}
\setlength{\emergencystretch}{2em}
\multlinegap0pt
\usepackage{bm}
\usepackage{bbm}
\usepackage{dsfont}
\usepackage{xspace}
\usepackage{cancel}
\usepackage{mathtools}
\usepackage{xcolor}
\usepackage{amsthm}
\makeatletter
\@ifundefined{thm}{\newtheorem{thm}{Theorem}}{}
\@ifundefined{lemma}{\newtheorem{lemma}[thm]{Lemma}}{}
\makeatother
\title[Non-commutative Intermediate Factor theorem associated with ]{Non-commutative Intermediate Factor theorem associated with $W^*$-dynamics of product groups}
\author{Tattwamasi Amrutam, Yongle Jiang,, Shuoxing Zhou}
\date{Source: arXiv:2508.18978v1 (CC BY 4.0)}
\begin{document}
\maketitle

\noindent\emph{Source and licence.} Non-commutative Intermediate Factor theorem associated with $W^*$-dynamics of product groups. Version arXiv:2508.18978v1 (CC BY 4.0), distributed under the Creative Commons Attribution 4.0 International licence, as recorded in the arXiv licence metadata for this exact version and shown on the abstract page https://arxiv.org/abs/2508.18978v1. Full rights notice: licenses/arxiv-2508.18978-CC-BY-4.0.txt.\par\smallskip

\noindent\textit{Editorial scope.} Counted unit: the Thm whose source label is \texttt{thm:genofBaderShalom} in the article (source0.tex lines 781--787, source label \texttt{thm:genofBaderShalom}), delivered together with the article's own complete original proof of it (source0.tex lines 788--806, one \texttt{proof} environment of 19 lines). No mathematics is written, repaired or re-derived: statement and proof are the article's own text line for line. Required context retained uncounted: lem:factor (lines 413--417); thm:vnsplitting (lines 427--433); lem:ergodic (lines 773--776). Every definition, display and label of those lines is retained unchanged. Omitted from this package and not redistributed: 1--412, 418--426, 434--772, 777--780, 807--853. Nothing inside a retained excerpt is omitted or altered: every retained body is delivered exactly as the article's lines are written, so no omission receipt of any kind accompanies this package. Citations: the retained excerpts carry no citation command at all, so no bibliography entry of the article is transcribed and no reference is redistributed. Reference adaptation: none was needed and none was made; every reference inside the delivered unit resolves inside it, so no unresolved reference or citation is produced. Printed slips kept literal: no printed word, symbol, hypothesis or number of the counted statement or of its proof was altered. Layout and class: the article's own class is \texttt{amsart} with packages including xcolor, ulem, dirtytalk, amsmath, amssymb, setspace, nicefrac, yhmath, amscd, eucal, srcltx, hyperref, amsrefs, tikz-cd; the wrapper is the lane's pinned \texttt{amsart} editorial wrapper at 10pt on A4, which restates the theorem-like environments the retained text needs and adds the article's own macro definitions that the retained text needs (none), with extra wrapper packages (bm, bbm, dsfont, xspace, cancel, mathtools, xcolor). The delivered numbering is the unit's own. Assets: no retained span contains a figure environment, a graphics-inclusion command, a PDF-page inclusion, a table or a TikZ picture; no image file is copied into this package, the package declares no asset dependency and no image file is redistributed with it. Licence: version arXiv:2508.18978v1 of this article is distributed under the Creative Commons Attribution 4.0 International licence, as recorded in the arXiv licence metadata for this exact version and shown on the abstract page https://arxiv.org/abs/2508.18978v1. A full rights notice is delivered at licenses/arxiv-2508.18978-CC-BY-4.0.txt. Characters: every retained excerpt is pure ASCII, so the delivered engine, XeLaTeX with its default Latin Modern text font, reports no missing character.
\begin{lemma}
\label{lem:factor}
Let $\Gamma$ be an ICC group. Assume that $(\mathcal{M},\tau)$ is a $(\Gamma,\mu)$-ergodic space in the sense that $\tau$ is $\mu$-stationary and $\text{Supp}(\mu)$ generates $\Gamma$. Then the action of $
\Gamma$ on the crossed product $\mathcal{M}\rtimes\Gamma$ by conjugation is also ergodic.    
\end{lemma}

\begin{thm}
\label{thm:vnsplitting}
Let $G$ be a countable ICC group. Let $\mathcal{M}$ be a $G$-von Neumann algebra. Let $H\le G$ be such that $(\mathcal{M},\tau)$ is a $(H,\mu)$-ergodic space.  Let $\mathcal{N}$ be a $G$-von Neumann algebra such that $H\curvearrowright \mathcal{N}$ is trivial. Then, every $G$-invariant von Neumann subalgebra $\mathcal{M}\bar{\otimes} \mathbb{C}\subset\mathcal{C}\subset \mathcal{M}\bar{\otimes} \mathcal{N}$ splits. 
\begin{proof}
  This is a direct corollary of the Ge-Kadison splitting theorem. To see this, set $\mathcal{Q}=\mathcal{M}\rtimes H$. It follows from Lemma~\ref{lem:factor} that $Q$ is a factor. Thus, by considering the inclusion of von Neumann algebras: $\mathcal{Q}\subseteq \mathcal{C}\rtimes H\subseteq (\mathcal{M}\bar{\otimes} \mathcal{N})\rtimes H=\mathcal{Q}\bar{\otimes} \mathcal{N}$. Hence $\mathcal{C}\rtimes H=\mathcal{Q}\bar{\otimes} \mathcal{P}=(\mathcal{M}\bar{\otimes} \mathcal{P})\rtimes H$ for some $\mathcal{P}\subseteq \mathcal{N}$. Then $\mathcal{C}=\mathcal{M}\bar{\otimes} \mathcal{P}$.
\end{proof}
\end{thm}

\begin{lemma}
\label{lem:ergodic}
 Let $(\mathcal{N},\tau)$ be a trace-preserving ergodic $G$-space. Let $(B,\nu_B)$ be the associated $(G,\mu)$-Poisson boundary. Then, $(\mathcal{N},\tau)\bar{\otimes}L^{\infty}(B,\nu_B)$ is also an ergodic $G$-space.
\end{lemma}

\begin{thm}
\label{thm:genofBaderShalom}
Let $G=G_1\times G_2$, $\mu=\mu_1\times\mu_2$. Assume that $\mu_i$'s are generating and $G_i$ is ICC Let $(\mathcal{N},\tau)$ be a trace-preserving $G_i$-space which is ergodic for each $G_i$. Let $(B,\nu_B)$ be a $(G_2,\mu_2)$-USB. Note that $G_1\curvearrowright (B_2,\nu_{B_2})$ is trivial. Consider an intermediate algebra $\mathcal{M}$ with 
\[(\mathcal{N},\tau)\subset\mathcal{M}\subset(\mathcal{N},\tau)\bar{\otimes}L^{\infty}(B,\nu_B)\]
Then,
$\mathcal{M}$ splits. 
\end{thm}

\begin{proof}
Let $\mathcal{N}_1=\left\{\mathcal{M},L^{\infty}(B_2,\nu_{B_2})\right\}''$. 
Since $G_2\curvearrowright (B_1,\nu_{B_1})$ is trivial, we see that

\[\left(L^{\infty}(B,\nu_{B})\bar{\otimes}(\mathcal{N},\tau)\right)\rtimes G_2=\left(\left((\mathcal{N},\tau)\otimes L^{\infty}(B_2,\nu_{B_2})\right)\rtimes G_2\right)\bar{\otimes} L^{\infty}(B_1,\nu_{B_1}).\]
We now see that 
\[\left((\mathcal{N},\tau)\otimes L^{\infty}(B_2,\nu_{B_2})\right)\rtimes G_2\subset\mathcal{N}_1\rtimes G_2\subset\left(L^{\infty}(B,\nu_B)\bar{\otimes}L^{\infty}(X_1,\xi_1)\right)\rtimes G_2.\]
This is equivalent to looking at
\begin{align*}&\left((\mathcal{N},\tau)\otimes L^{\infty}(B_2,\nu_{B_2})\right)\rtimes G_2\\&\subset\mathcal{N}_1\rtimes G_2\\&\subset\left(\left((\mathcal{N},\tau)\otimes L^{\infty}(B_2,\nu_{B_2})\right)\rtimes G_2\right)\bar{\otimes} L^{\infty}(B_1,\nu_{B_1})\end{align*}
For $H=G_2$, using Theorem~\ref{thm:vnsplitting} and Lemma~\ref{lem:ergodic}, we see that
\[\mathcal{N}_1\rtimes G_2=\left(\left((\mathcal{N},\tau)\otimes L^{\infty}(B_2,\nu_{B_2})\right)\rtimes G_2\right)\bar{\otimes} L^{\infty}(\tilde{B}_1,\nu_{\tilde{B}_1})\]
Since $G_2\curvearrowright(\tilde{B}_1,\nu_{\tilde{B}_1})$ is trivial, applying the canonical conditional expectation, we get that
\[\mathcal{N}_1=\left((\mathcal{N},\tau)\otimes L^{\infty}(B_2,\nu_{B_2})\right)\bar{\otimes}L^{\infty}(\tilde{B}_1,\nu_{\tilde{B}_1}).\]
\textcolor{teal}{We claim that $L^{\infty}(\tilde{B}_1,\nu_{\tilde{B}_1})\subset \mathcal{M}$}. Once this is established, we will be done. Indeed, let us assume that the claim holds and finish the proof. If the claim holds, we obtain that
\[(\mathcal{N},\tau)\bar{\otimes} L^{\infty}(\tilde{B}_1,\nu_{\tilde{B}_1})\subset\mathcal{M}\subset\left((\mathcal{N},\tau)\bar{\otimes} L^{\infty}(\tilde{B}_1,\nu_{\tilde{B}_1})\right)\bar{\otimes} L^{\infty}(B_2,\nu_{B_2}).\]
Using Theorem~\ref{thm:vnsplitting} along with Lemma~\ref{lem:ergodic} for $H=G_1$, we obtain that
\[\mathcal{M}=(\mathcal{N},\tau)\bar{\otimes} L^{\infty}(\tilde{B}_1,\nu_{\tilde{B}_1})\bar{\otimes} L^{\infty}(\tilde{B}_2,\nu_{\tilde{B}_2}).\]
\textcolor{red}{\textit{Proof of Claim}}???
\end{proof}

\end{document}
